\documentclass[10pt,a4paper]{amsart}
\usepackage{amsmath,amssymb,amsthm,graphicx}
\usepackage{verbatim}
\newenvironment{reference}{\comment}{\endcomment}
%\newenvironment{reference}{}{}
\newenvironment{slogan}{\comment}{\endcomment}
\newenvironment{history}{\comment}{\endcomment}

% For commutative diagrams we use Xy-pic
\usepackage[all]{xy}

% We use 2cell for 2-commutative diagrams.
\xyoption{2cell}
% Theorem environments.
%
\theoremstyle{plain}
\newtheorem{theorem}[subsection]{Theorem}
\newtheorem{proposition}[subsection]{Proposition}
\newtheorem{lemma}[subsection]{Lemma}

\theoremstyle{definition}
\newtheorem{definition}[subsection]{Definition}
\newtheorem{example}[subsection]{Example}
\newtheorem{exercise}[subsection]{Exercise}
\newtheorem{situation}[subsection]{Situation}

\theoremstyle{remark}
\newtheorem{remark}[subsection]{Remark}
\newtheorem{remarks}[subsection]{Remarks}

\numberwithin{equation}{subsection}

% Macros
%
\def\lim{\mathop{\mathrm{lim}}\nolimits}
\def\colim{\mathop{\mathrm{colim}}\nolimits}
\def\Spec{\mathop{\mathrm{Spec}}}
\def\Hom{\mathop{\mathrm{Hom}}\nolimits}
\def\Ext{\mathop{\mathrm{Ext}}\nolimits}
\def\SheafHom{\mathop{\mathcal{H}\!\mathit{om}}\nolimits}
\def\SheafExt{\mathop{\mathcal{E}\!\mathit{xt}}\nolimits}
\def\Sch{\mathit{Sch}}
\def\Mor{\mathop{\mathrm{Mor}}\nolimits}
\def\Ob{\mathop{\mathrm{Ob}}\nolimits}
\def\Sh{\mathop{\mathit{Sh}}\nolimits}
\def\NL{\mathop{N\!L}\nolimits}
\def\CH{\mathop{\mathrm{CH}}\nolimits}
\def\proetale{{pro\text{-}\acute{e}tale}}
\def\etale{{\acute{e}tale}}
\def\QCoh{\mathit{QCoh}}
\def\Ker{\mathop{\mathrm{Ker}}}
\def\Im{\mathop{\mathrm{Im}}}
\def\Coker{\mathop{\mathrm{Coker}}}
\def\Coim{\mathop{\mathrm{Coim}}}

% Boxtimes
%
\DeclareMathSymbol{\boxtimes}{\mathbin}{AMSa}{"02}

%
% Macros for moduli stacks/spaces
%
\def\QCohstack{\mathcal{QC}\!\mathit{oh}}
\def\Cohstack{\mathcal{C}\!\mathit{oh}}
\def\Spacesstack{\mathcal{S}\!\mathit{paces}}
\def\Quotfunctor{\mathrm{Quot}}
\def\Hilbfunctor{\mathrm{Hilb}}
\def\Curvesstack{\mathcal{C}\!\mathit{urves}}
\def\Polarizedstack{\mathcal{P}\!\mathit{olarized}}
\def\Complexesstack{\mathcal{C}\!\mathit{omplexes}}
% \Pic is the operator that assigns to X its picard group, usage \Pic(X)
% \Picardstack_{X/B} denotes the Picard stack of X over B
% \Picardfunctor_{X/B} denotes the Picard functor of X over B
\def\Pic{\mathop{\mathrm{Pic}}\nolimits}
\def\Picardstack{\mathcal{P}\!\mathit{ic}}
\def\Picardfunctor{\mathrm{Pic}}
\def\Deformationcategory{\mathcal{D}\!\mathit{ef}}
\usepackage{cleveref}
\crefname{originalreference}{original reference}{original references}
\Crefname{originalreference}{Original reference}{Original references}
\makeatletter
\newlabel{lemma-affine-coherator}{{08D9}{source}{}{sourcepointer-1}{}}
\newlabel{lemma-affine-coherator@cref}{{[originalreference][0][]08D9}{[1][0][]source}{}{}{}}
\makeatother
\crefname{originalreference}{original reference}{original references}
\Crefname{originalreference}{Original reference}{Original references}
\makeatletter
\newlabel{lemma-better-coherator}{{0CR0}{source}{}{sourcepointer-2}{}}
\newlabel{lemma-better-coherator@cref}{{[originalreference][0][]0CR0}{[1][0][]source}{}{}{}}
\makeatother
\crefname{originalreference}{original reference}{original references}
\Crefname{originalreference}{Original reference}{Original references}
\makeatletter
\newlabel{lemma-induction-principle}{{08DR}{source}{}{sourcepointer-3}{}}
\newlabel{lemma-induction-principle@cref}{{[originalreference][0][]08DR}{[1][0][]source}{}{}{}}
\makeatother
\crefname{originalreference}{original reference}{original references}
\Crefname{originalreference}{Original reference}{Original references}
\makeatletter
\newlabel{lemma-quasi-coherence-direct-image}{{08D5}{source}{}{sourcepointer-4}{}}
\newlabel{lemma-quasi-coherence-direct-image@cref}{{[originalreference][0][]08D5}{[1][0][]source}{}{}{}}
\makeatother
\crefname{originalreference}{original reference}{original references}
\Crefname{originalreference}{Original reference}{Original references}
\makeatletter
\newlabel{remark-explain-consequence}{{0CR2}{source}{}{sourcepointer-5}{}}
\newlabel{remark-explain-consequence@cref}{{[originalreference][0][]0CR2}{[1][0][]source}{}{}{}}
\makeatother
\title{The coherator DQ shifts boundedness by a uniform amount depending only on X}
\author{The Stacks Project Authors}
\date{stacks-project:spaces-perfect.tex:local-snapshot-20261012; GNU FDL 1.2 or later}
\begin{document}
\maketitle
\noindent\textit{Editorial note.} All delivered mathematical text is the author's own, copied from the pinned Stacks Project chapter file spaces-perfect.tex; no statement or proof step was rewritten, reordered or removed.\par\medskip
\noindent\textit{Editorial note.} The counted proof is the author's complete argument: the proof runs the induction principle with the affine case as base, controls the three pieces of the distinguished square by their individual bounds, and absorbs them into a single N via the finite cohomological dimension of the Rj\_* functors\par\medskip
\noindent\textit{Editorial note.} The label lemma-boundedness-better-coherator names the counted statement, and the proof environment immediately following it in the source is delivered verbatim as its proof.\par\medskip
\noindent\textit{Editorial note.} The counted statement is carried by the author's own lemma environment in the chapter source.\par\medskip
\noindent\textit{Allowed standard prerequisites (editorial).}\begin{itemize}
\item the coherator DQ\_X of lemma-better-coherator and its affine description RQ\_X = DQ\_X as the quasi-coherent sheaf attached to RGamma(X, K) (lemma-affine-coherator of the same chapter)
\item the induction principle for properties of quasi-compact quasi-separated algebraic spaces (lemma-induction-principle of the same chapter)
\item elementary distinguished squares (U subset W, f : V -> W) with quasi-compact quasi-separated W, quasi-compact open U and affine V, and the Mayer-Vietoris distinguished triangle of remark-explain-consequence
\item finite cohomological dimension on D\_QCoh of the functors Rj\_* for the thickenings j of the square (lemma-quasi-coherence-direct-image of the same chapter)
\end{itemize}\medskip
\ifdefined\hypertarget\hypertarget{sourcepointer-1}{}\fi
\noindent\textit{Source reference.} Original reference 08D9: sources/stacks-spaces-perfect/pinned-original-source.tex, source lines 2441--2441. Prerequisite pointer invoked at proof line(s) 4807; it names a standard fact of the same chapter that the counted argument uses, and every step of the counted proof built on it is retained verbatim in the delivered body.\par\medskip
\ifdefined\hypertarget\hypertarget{sourcepointer-2}{}\fi
\noindent\textit{Source reference.} Original reference 0CR0: sources/stacks-spaces-perfect/pinned-original-source.tex, source lines 4654--4654. Prerequisite pointer invoked at proof line(s) 4790; it names a standard fact of the same chapter that the counted argument uses, and every step of the counted proof built on it is retained verbatim in the delivered body.\par\medskip
\ifdefined\hypertarget\hypertarget{sourcepointer-3}{}\fi
\noindent\textit{Source reference.} Original reference 08DR: sources/stacks-spaces-perfect/pinned-original-source.tex, source lines 1563--1563. Prerequisite pointer invoked at proof line(s) 4801; it names a standard fact of the same chapter that the counted argument uses, and every step of the counted proof built on it is retained verbatim in the delivered body.\par\medskip
\ifdefined\hypertarget\hypertarget{sourcepointer-4}{}\fi
\noindent\textit{Source reference.} Original reference 08D5: sources/stacks-spaces-perfect/pinned-original-source.tex, source lines 790--790. Prerequisite pointer invoked at proof line(s) 4823; it names a standard fact of the same chapter that the counted argument uses, and every step of the counted proof built on it is retained verbatim in the delivered body.\par\medskip
\ifdefined\hypertarget\hypertarget{sourcepointer-5}{}\fi
\noindent\textit{Source reference.} Original reference 0CR2: sources/stacks-spaces-perfect/pinned-original-source.tex, source lines 4769--4769. Prerequisite pointer invoked at proof line(s) 4829; it names a standard fact of the same chapter that the counted argument uses, and every step of the counted proof built on it is retained verbatim in the delivered body.\par\medskip
\begin{lemma}
\label{lemma-boundedness-better-coherator}
Let $S$ be a scheme.
Let $X$ be a quasi-compact and quasi-separated algebraic space over $S$.
The functor $DQ_X$ of Lemma \ref{lemma-better-coherator}
has the following boundedness property:
there exists an integer $N = N(X)$ such that, if
$K$ in $D(\mathcal{O}_X)$ with
$H^i(U, K) = 0$ for $U$ affine \'etale over $X$ and $i \not \in [a, b]$, then
the cohomology sheaves $H^i(DQ_X(K))$ are zero for
$i \not \in [a, b + N]$.
\end{lemma}

\begin{proof}
We will prove this using the induction principle of
Lemma \ref{lemma-induction-principle}.

\medskip\noindent
If $X$ is affine, then the lemma is true with $N = 0$ because then
$RQ_X = DQ_X$ is given by taking the complex of
quasi-coherent sheaves associated to $R\Gamma(X, K)$.
See Lemma \ref{lemma-affine-coherator}.

\medskip\noindent
Let $(U \subset W, f : V \to W)$ be an elementary distinguished square
with $W$ quasi-compact and quasi-separated, $U \subset W$
quasi-compact open, $V$ affine such that
the lemma holds for $U$, $V$, and $U \times_W V$.
Say with integers $N(U)$, $N(V)$, and $N(U \times_W V)$.
Now suppose $K$ is in $D(\mathcal{O}_X)$ with
$H^i(W, K) = 0$ for all affine $W$ \'etale over $X$ and all $i \not \in [a, b]$.
Then $K|_U$, $K|_V$, $K|_{U \times_W V}$ have the same property.
Hence we see that $RQ_U(K|_U)$ and $RQ_V(K|_V)$ and
$RQ_{U \cap V}(K|_{U \times_W V})$ have vanishing cohomology
sheaves outside the interval $[a, b + \max(N(U), N(V), N(U \times_W V))$.
Since the functors $Rj_{U, *}$, $Rj_{V, *}$, $Rj_{U \times_W V, *}$
have finite cohomological dimension on $D_\QCoh$ by
Lemma \ref{lemma-quasi-coherence-direct-image}
we see that there exists an $N$ such that
$Rj_{U, *}DQ_U(K|_U)$, $Rj_{V, *}DQ_V(K|_V)$, and
$Rj_{U \cap V, *}DQ_{U \times_W V}(K|_{U \times_W V})$ have vanishing
cohomology sheaves outside the interval $[a, b + N]$.
Then finally we conclude by the distinguished triangle
of Remark \ref{remark-explain-consequence}.
\end{proof}

\end{document}
